% =====================================================================
%  Thème 4 — Produit scalaire et orthogonalité — FACILE — Exercices
% =====================================================================
\documentclass[11pt,a4paper]{article}
\usepackage[utf8]{inputenc}
\usepackage[T1]{fontenc}
\usepackage{lmodern}
\usepackage[french]{babel}
\usepackage{amsmath,amssymb,mathtools}
\usepackage{xcolor}
\usepackage{enumitem}
\usepackage{fancyhdr}
\usepackage[a4paper,margin=2.1cm,headheight=26pt,headsep=10pt]{geometry}

\definecolor{primary}{RGB}{222,70,80}
\definecolor{primarydark}{RGB}{150,30,42}

\pagestyle{fancy}\fancyhf{}
\fancyhead[L]{\small\textcolor{primarydark}{\textbf{Fiche 2} — Calcul vectoriel}}
\fancyhead[R]{\small\textcolor{primarydark}{\textsc{Thème 4 — Facile — Exercices}}}
\fancyfoot[C]{\small\thepage}
\renewcommand{\headrulewidth}{0.8pt}

\newcommand{\vc}[1]{\overrightarrow{#1}}
\providecommand{\norm}[1]{\left\lVert #1\right\rVert}
\newcommand{\exo}[1]{\medskip\noindent{\color{primary}\bfseries\large Exercice #1.}\par\nopagebreak\smallskip}

\begin{document}
\begin{center}
{\LARGE\bfseries\color{primarydark}Produit scalaire et orthogonalité — Niveau Facile}\\[2pt]
{\color{primary}\rule{0.55\linewidth}{1.1pt}}
\end{center}
\vspace{1mm}
{\small\itshape Objectif : calculer un produit scalaire avec $xx'+yy'$, tester l'orthogonalité et
utiliser le carré scalaire.}
\vspace{2mm}

\exo{1} \textbf{Calculer un produit scalaire.}
Calculer $\vec u\cdot\vec v$ pour chaque couple.
\begin{enumerate}[label=\textbf{\alph*)},leftmargin=2em,itemsep=2pt]
  \item $\vec u\,(3;2)$, $\vec v\,(1;4)$.
  \item $\vec u\,(5;-1)$, $\vec v\,(2;3)$.
  \item $\vec u\,(-2;4)$, $\vec v\,(3;-1)$.
  \item $\vec u\,(2;0)$, $\vec v\,(0;5)$.
\end{enumerate}

\exo{2} \textbf{Orthogonaux ou non ?}
Pour chaque couple, calculer le produit scalaire et dire si les vecteurs sont orthogonaux.
\begin{enumerate}[label=\textbf{\alph*)},leftmargin=2em,itemsep=2pt]
  \item $\vec u\,(3;2)$, $\vec v\,(-2;3)$.
  \item $\vec u\,(1;4)$, $\vec v\,(2;3)$.
  \item $\vec u\,(6;-3)$, $\vec v\,(1;2)$.
\end{enumerate}

\exo{3} \textbf{Carré scalaire.}
Calculer $\vec u\cdot\vec u$, puis en déduire $\norm{\vec u}$.
\begin{enumerate}[label=\textbf{\alph*)},leftmargin=2em,itemsep=2pt]
  \item $\vec u\,(3;4)$.
  \item $\vec u\,(1;-2)$.
\end{enumerate}

\exo{4} \textbf{Signe du produit scalaire et nature de l'angle.}
Sans calculer l'angle, dire s'il est aigu, droit ou obtus.
\begin{enumerate}[label=\textbf{\alph*)},leftmargin=2em,itemsep=2pt]
  \item $\vec u\,(2;1)$, $\vec v\,(3;4)$.
  \item $\vec u\,(2;1)$, $\vec v\,(-3;1)$.
  \item $\vec u\,(1;2)$, $\vec v\,(2;-1)$.
\end{enumerate}

\exo{5} \textbf{Droites perpendiculaires (vecteurs directeurs donnés).}
Deux droites ont pour vecteurs directeurs les vecteurs indiqués. Dire si elles sont perpendiculaires.
\begin{enumerate}[label=\textbf{\alph*)},leftmargin=2em,itemsep=2pt]
  \item $\vec u\,(1;2)$ et $\vec v\,(2;-1)$.
  \item $\vec u\,(3;1)$ et $\vec v\,(1;3)$.
\end{enumerate}

\end{document}
